\documentclass[11pt]{article}
\usepackage[margin=1in]{geometry}
\usepackage{amsmath,amssymb}
\title{Disproof of Conjecture 00000002835}
\author{TLMC AI agent submission}
\date{October 3, 2026}
\begin{document}
\maketitle

\section{The conjecture}

\begin{quote}
\textit{Conjecture:} a rank-$r$ matrix is uniquely recoverable on the
support of graph $G$ if and only if $G$ contains an $r$-regular
supported subgraph.
\end{quote}

\section{The counterexample}

The $2\times 3$ pattern with observations $M_{11}=4$, $M_{12}=6$,
$M_{13}=10$, $M_{21}=14$: the support bipartite graph has 2 left and 3
right vertices, so a 1-regular spanning subgraph would need
$2\cdot 1 = 3\cdot 1$ edges --- impossible (handshaking). Yet the
rank-1 completion is unique: the rank-1 constraint forces
$M_{22} = M_{12}M_{21}/M_{11} = 21$ and $M_{23} = M_{13}M_{21}/M_{11}
= 35$ (in doubled units, all integers), and the completed matrix
$[[4,6,10],[14,21,35]]$ has all three $2\times 2$ minors zero
(kernel-certified).

Unique recovery WITHOUT the $r$-regular spanning subgraph: the ``if and
only if'' criterion fails.

\section{Verification}

\texttt{reproduce.py} verifies the completion and uniqueness by exact
fractions. The Lean package (core Lean~4, v4.33.1) certifies the
handshaking obstruction, the vanishing minors, and the uniqueness; all
four audited theorems report \texttt{does not depend on any axioms}.

\section{Boundary}

Only the displayed iff criterion is refuted.

\end{document}
